\begin{ExerciseBlock}
\ExerciseHeading{03}{Zeit- und Ortsabhängigkeit einer ebenen Welle}{wave_propagation_harmonic_travelling_wave}
Eine Quelle bei $x=0$ regt ab $t=0$ eine harmonische ebene Transversalwelle an, die sich mit \qty{3}{\metre\per\second} in positiver $x$-Richtung ausbreitet. An der Quelle gilt $y(0,0)=0$. Die Auslenkung steigt innerhalb einer Sekunde erstmals auf $y_0=\qty{10}{\centi\metre}$.
\begin{enumerate}
\item Bestimmen Sie die Periodendauer $T$, die Kreisfrequenz $\omega$ und die Frequenz $f$.
\item Welche Wellenlänge $\lambda$ und Wellenzahl $k$ hat die Welle?
\item Geben Sie die spezielle Lösung $y(x,t)$ hinter der Wellenfront an.
\item Nach welcher Zeit erreicht die Wellenfront einen Punkt in \qty{120}{\metre} Entfernung?
\item Welche Auslenkung herrscht bei $t=\qty{50}{\second}$ und $x=\qty{50}{\metre}$?
\end{enumerate}
\end{ExerciseBlock}
\UseExerciseSolution{wave_propagation_harmonic_travelling_wave}
